% GENERATED FILE. Edit canonical lessons, then run npm run generate:books.
% canonical-source-hash: fnv1a64:444b93e801f10cf8
% canonical-lessons: RU-C290-moshennik, RU-C290-obmanyvat, RU-C290-ugonyat, RU-C290-pokhishchat, RU-C290-pryatatsya

\chapter{Tricks and hiding}
\label{ch:ru-tricks-and-hiding}
\hlchaptermodality{290}

\begin{chapteropening}
\textbf{By the end of this chapter:} \emph{I can talk about a swindler, cheating, stealing a car, kidnapping and hiding.}

Complete the last lesson of chapter 290: I can talk about a swindler, cheating, stealing a car, kidnapping and hiding.
\end{chapteropening}

\section[\texorpdfstring{\ru{мошенник} (moshénnik)}{мошенник (moshénnik)}]{\ru{мошенник} (moshénnik) — a swindler, a fraudster}

\label{lesson:RU-C290-moshennik}



\begin{quote}
Before the new one: say the Russian for an attack, then the Russian for to attack.
\end{quote}

\subsection*{You'll want to know: \ru{мошенник}}
\textbf{\ru{мошенник}} (\emph{moshénnik}) — \textquotedblleft{}a swindler, a fraudster\textquotedblright{}.

\begin{quote}
\ru{Мошенник} \ru{позвонил} \ru{бабушке}.
\end{quote}

\begin{quote}
\emph{Moshénnik pozvoníl bábushke.}
\end{quote}

\begin{quote}
A swindler phoned Grandma.
\end{quote}

\emph{\ru{Мошенничество}} — fraud.

\begin{cousinweb}
\emph{\ru{мошенник}} is usually traced to the old word \emph{\ru{мошна}} \textquotedblleft{}a purse\textquotedblright{}, so first \textquotedblleft{}a cutpurse\textquotedblright{}; EtymDB lists no parent.
\end{cousinweb}

\begin{grammarlens}[title={What you've built}]
One more word for telling, explaining and describing at B1.
\end{grammarlens}

\subsection*{Your turn}
\begin{itemize}
\raggedright
  \item \emph{Say it:} \emph{moshénnik}
  \item \emph{Say it:} \emph{moshénnik}, once more
  \item \emph{Say it:} a sentence of your own with \emph{moshénnik}
  \item \emph{Recall:} say the Russian for an attack, then the Russian for to attack, then say \emph{moshénnik} again
\end{itemize}

\begin{tcolorbox}[breakable,colback=teal!4,colframe=teal!35!black,title={Before you move on}]
What does \ru{мошенник} mean? (\textquotedblleft{}A swindler, a fraudster\textquotedblright{}.) Say it once more.
\end{tcolorbox}
\section[\texorpdfstring{\ru{обманывать} (obmányvat)}{обманывать (obmányvat)}]{\ru{обманывать} (obmányvat) — to deceive, to cheat}

\label{lesson:RU-C290-obmanyvat}



\begin{quote}
Before the new one: say the Russian for to attack, then the Russian for a swindler.
\end{quote}

\subsection*{You'll want to know: \ru{обманывать}}
\textbf{\ru{обманывать}} (\emph{obmányvat}) — \textquotedblleft{}to deceive, to cheat\textquotedblright{}.

\begin{quote}
\ru{Он} \ru{никогда} \ru{не} \ru{обманывает} \ru{друзей}.
\end{quote}

\begin{quote}
\emph{On nikogdá ne obmányvayet druzéy.}
\end{quote}

\begin{quote}
He never deceives his friends.
\end{quote}

Imperfective; perfective \emph{\ru{обмануть}} (\emph{obmanút}). \emph{\ru{Обман}} — a deceit, a trick.

\begin{cousinweb}
\emph{\ru{обманывать}} is built on \emph{\ru{об}-} and \emph{\ru{манить}} \textquotedblleft{}to lure, to beckon\textquotedblright{}, which is inherited from Proto-Slavic \emph{maniti} \textquotedblleft{}to deceive\textquotedblright{} (EtymDB).
\end{cousinweb}

\begin{grammarlens}[title={What you've built}]
One more word for telling, explaining and describing at B1.
\end{grammarlens}

\subsection*{Your turn}
\begin{itemize}
\raggedright
  \item \emph{Say it:} \emph{obmányvat}
  \item \emph{Say it:} \emph{obmányvat}, once more
  \item \emph{Say it:} a sentence of your own with \emph{obmányvat}
  \item \emph{Recall:} say the Russian for to attack, then the Russian for a swindler, then say \emph{obmányvat} again
\end{itemize}

\begin{tcolorbox}[breakable,colback=teal!4,colframe=teal!35!black,title={Before you move on}]
What does \ru{обманывать} mean? (\textquotedblleft{}To deceive, to cheat\textquotedblright{}.) Say it once more.
\end{tcolorbox}
\section[\texorpdfstring{\ru{угонять} (ugonyát)}{угонять (ugonyát)}]{\ru{угонять} (ugonyát) — to hijack, to steal (a vehicle)}

\label{lesson:RU-C290-ugonyat}



\begin{quote}
Before the new one: say the Russian for a swindler, then the Russian for to deceive.
\end{quote}

\subsection*{You'll want to know: \ru{угонять}}
\textbf{\ru{угонять}} (\emph{ugonyát}) — \textquotedblleft{}to hijack, to steal (a vehicle)\textquotedblright{}.

\begin{quote}
\ru{Ночью} \ru{у} \ru{соседа} \ru{угнали} \ru{машину}.
\end{quote}

\begin{quote}
\emph{Nóchyu u soséda ugnáli mashínu.}
\end{quote}

\begin{quote}
During the night the neighbour's car was stolen.
\end{quote}

Imperfective; perfective \emph{\ru{угнать}} (\emph{ugnát}), \emph{\ru{угоню}}. For a car or a plane; \emph{\ru{угон}} — a car theft.

\begin{cousinweb}
\emph{\ru{угонять}} is built on \emph{\ru{у}-} \textquotedblleft{}away\textquotedblright{} and \emph{\ru{гнать}} \textquotedblleft{}to drive\textquotedblright{}, which is inherited from Proto-Slavic \emph{gŭnati}, from PIE \emph{g\textsuperscript{w}\textsuperscript{h}en-} \textquotedblleft{}to strike\textquotedblright{} (EtymDB): to drive away.
\end{cousinweb}

\begin{grammarlens}[title={What you've built}]
One more word for telling, explaining and describing at B1.
\end{grammarlens}

\subsection*{Your turn}
\begin{itemize}
\raggedright
  \item \emph{Say it:} \emph{ugonyát}
  \item \emph{Say it:} \emph{ugonyát}, once more
  \item \emph{Say it:} a sentence of your own with \emph{ugonyát}
  \item \emph{Recall:} say the Russian for a swindler, then the Russian for to deceive, then say \emph{ugonyát} again
\end{itemize}

\begin{tcolorbox}[breakable,colback=teal!4,colframe=teal!35!black,title={Before you move on}]
What does \ru{угонять} mean? (\textquotedblleft{}To hijack, to steal (a vehicle)\textquotedblright{}.) Say it once more.
\end{tcolorbox}
\section[\texorpdfstring{\ru{похищать} (pokhishchát)}{похищать (pokhishchát)}]{\ru{похищать} (pokhishchát) — to kidnap, to abduct; to steal}

\label{lesson:RU-C290-pokhishchat}



\begin{quote}
Before the new one: say the Russian for to deceive, then the Russian for to hijack.
\end{quote}

\subsection*{You'll want to know: \ru{похищать}}
\textbf{\ru{похищать}} (\emph{pokhishchát}) — \textquotedblleft{}to kidnap, to abduct; to steal\textquotedblright{}.

\begin{quote}
\ru{В} \ru{фильме} \ru{похитили} \ru{дочь} \ru{богатого} \ru{человека}.
\end{quote}

\begin{quote}
\emph{V fílme pokhítili doch bogátovo chelovéka.}
\end{quote}

\begin{quote}
In the film the daughter of a rich man was kidnapped.
\end{quote}

Imperfective; perfective \emph{\ru{похитить}} (\emph{pokhítit}), \emph{\ru{я} \ru{похищу}}. Bookish.

\begin{cousinweb}
\emph{\ru{похищать}} is built on \emph{\ru{по}-} and an old verb \emph{\ru{хитить}} \textquotedblleft{}to seize\textquotedblright{}, the root of \emph{\ru{хищник}} \textquotedblleft{}a predator\textquotedblright{}; EtymDB lists no parent.
\end{cousinweb}

\begin{grammarlens}[title={What you've built}]
One more word for telling, explaining and describing at B1.
\end{grammarlens}

\subsection*{Your turn}
\begin{itemize}
\raggedright
  \item \emph{Say it:} \emph{pokhishchát}
  \item \emph{Say it:} \emph{pokhishchát}, once more
  \item \emph{Say it:} a sentence of your own with \emph{pokhishchát}
  \item \emph{Recall:} say the Russian for to deceive, then the Russian for to hijack, then say \emph{pokhishchát} again
\end{itemize}

\begin{tcolorbox}[breakable,colback=teal!4,colframe=teal!35!black,title={Before you move on}]
What does \ru{похищать} mean? (\textquotedblleft{}To kidnap, to abduct; to steal\textquotedblright{}.) Say it once more.
\end{tcolorbox}
\section[\texorpdfstring{\ru{прятаться} (pryátatsya)}{прятаться (pryátatsya)}]{\ru{прятаться} (pryátatsya) — to hide (oneself)}

\label{lesson:RU-C290-pryatatsya}



\begin{quote}
Before the new one: say the Russian for to hijack, then the Russian for to kidnap.
\end{quote}

\subsection*{You'll want to know: \ru{прятаться}}
\textbf{\ru{прятаться}} (\emph{pryátatsya}) — \textquotedblleft{}to hide (oneself)\textquotedblright{}.

\begin{quote}
\ru{Кошка} \ru{прячется} \ru{под} \ru{кроватью}.
\end{quote}

\begin{quote}
\emph{Kóshka pryáchetsya pod krovátyu.}
\end{quote}

\begin{quote}
The cat is hiding under the bed.
\end{quote}

Imperfective; perfective \emph{\ru{спрятаться}} (\emph{spryátatsya}). Present \emph{\ru{прячусь}}, \emph{\ru{прячешься}}. Without \emph{-\ru{ся}}, \emph{\ru{прятать}} hides something else.

\begin{cousinweb}
\emph{\ru{прятаться}} is \emph{\ru{прятать}} \textquotedblleft{}to hide\textquotedblright{} with the reflexive \emph{-\ru{ся}} (EtymDB): to hide oneself. The origin of \emph{\ru{прятать}} is not settled.
\end{cousinweb}

\begin{grammarlens}[title={What you've built}]
One more word for telling, explaining and describing at B1.
\end{grammarlens}

\subsection*{Your turn}
\begin{itemize}
\raggedright
  \item \emph{Say it:} \emph{pryátatsya}
  \item \emph{Say it:} \emph{pryátatsya}, once more
  \item \emph{Say it:} a sentence of your own with \emph{pryátatsya}
  \item \emph{Recall:} say the Russian for to hijack, then the Russian for to kidnap, then say \emph{pryátatsya} again
\end{itemize}

\begin{tcolorbox}[breakable,colback=teal!4,colframe=teal!35!black,title={Before you move on}]
What does \ru{прятаться} mean? (\textquotedblleft{}To hide (oneself)\textquotedblright{}.) Say it once more.
\end{tcolorbox}
